\section{Model Discussion}

\subsection{Strengths}

\begin{enumerate}
    \item
    The model calculates how cooled air mixes with outdoor air before estimating the temperature in the waiting area. It also limits evaporative cooling by the outdoor wet-bulb temperature and accounts for the moisture added to the air.

    \item
    Standard Effective Temperature combines air temperature, humidity, airspeed, and radiant temperature into a common comfort measure. This allows the model to capture the benefit of a fan even when it leaves the air temperature unchanged.

    \item
    We report SET reduction per kWh separately from hot-discomfort loss and minimax regret.

    \item
    Minimax regret ranks equipment at each venue by its largest regret across three hot-weather scenarios. At Grove Hall, we tested how limits on purchase cost, electricity, and water changed the selection.

    \item
    The sensitivity sweeps locate intervals where the recommendation changes as individual inputs vary.
\end{enumerate}

\subsection{Weaknesses}

\begin{enumerate}
    \item
    The model estimates average conditions in each waiting area. This keeps the comparison consistent across venues, although cooling at a particular seat may differ from the zone average.

    \item
    We use the uncooled outdoor air temperature as a proxy for mean radiant temperature. Passenger-level measurements could assess how shade and hot surfaces affect this approximation.

    \item
    Each venue has three modeled weather cases. Different weather or resource limits could change the ranking. The sweeps consider individual inputs, so joint effects remain uncertain.
\end{enumerate}